% ===========================================================================
%  Z8440/1/2/4, Z84C40/1/2/3/4 Serial Input/Output Controller
%  Recreation of Zilog product specification PS018301-0602, pp. 99-122.
%  STYLE PROOF: p.99 complete + Figures 1 and 2.
% ===========================================================================
\documentclass[10pt,twoside]{article}
\usepackage[letterpaper,top=0.80in,bottom=0.80in,left=0.95in,right=0.95in,
            footskip=0.42in]{geometry}
\usepackage{z80sio}

\begin{document}
\setcounter{page}{99}

% ------------------------------- PAGE 99 -----------------------------------
\masthead{{\scshape Product Specification}}
         {Z8440/1/2/4,\\Z84C40/1/2/3/4}
         {{\scshape Serial Input/Output Controller}}

\dsrule
\dshead{FEATURES}

\begin{dscols}
\begin{featlist}
\item Two independent full-duplex channels, with separate control and status
      lines for modems or other devices.
\item Data rate in the x1 clock mode of 0 to 2.0M bits/second with a
      10 MHz clock.
\item NMOS version for cost sensitive performance solutions, CMOS version
      for the designs requiring low power consumption.
\item NMOS Z0844x04 -- 4 MHz, Z0844x06 -- 6.17 MHz (Where x is the designator
      for the bonding option; 0, 1, 2 or 4)
\item CMOS Z84C4x06 -- DC to 6/7 MHz, Z84C4x08 -- DC to 8 MHz, Z84C4x10 --
      DC to 10 MHz (Where x is the designator for the bonding option,
      0, 1, 2, or 4)
\item 6 MHz version supports 6.144 MHz CPU clock operation.
\end{featlist}

\columnbreak

\begin{featlist}
\item Asynchronous protocols: everything necessary for complete messages in
      5, 6, 7, or 8 bits/character. Includes variable stop bits and several
      clock-rate multipliers; break generation and detection; parity; overrun
      and framing error detection.
\item Synchronous protocols: everything necessary for complete bit- or
      byte-oriented messages in 5, 6, 7, or 8 bits/character, including IBM
      Bisync, SDLC, HDLC, CCITT-X.25 and others. Automatic CRC
      generation/checking, sync character and zero insertion/deletion, abort
      generation/detection, and flag insertion.
\item Receiver data registers quadruply buffered, transmitter registers doubly
      buffered.
\item Highly sophisticated and flexible daisy-chain interrupt vectoring for
      interrupts without external logic.
\end{featlist}
\end{dscols}

\dsrule
\dshead{GENERAL DESCRIPTION}

\begin{dscols}
The Z80 SIO (hereinafter referred to as the Z80 SIO or, SIO) Serial
Input/Output Controller is a dual-channel data communication interface with
extraordinary versatility and capability. Its basic functions as a
serial-to-parallel, parallel-to-serial converter/controller can be programmed
by a CPU for a broad range of serial communication applications.

The device supports all common asynchronous and synchronous protocols, byte-
or bit-oriented, and performs all of the functions traditionally done by
UARTs, USARTs, and synchronous communication controllers combined, plus
additional functions traditionally performed by the CPU. Moreover, it does
this on two fully-independent channels, with an exceptionally sophisticated
interrupt structure that allows very fast transfers.

Full interfacing is provided for CPU or DMA control. In addition to data
communication, the circuit can handle virtually all types of serial I/O with
fast, or slow, peripheral devices. While designed primarily as a member of the
Z80 family, its versatility makes it well suited to many other CPUs.

\textbf{The Z80 SIO uses a single +5V power supply and the standard Z80 family
single-phase clock. The SIO/0, SIO/1, and SIO/2 are packaged in a 40-pin DIP,
the SIO/4 is packaged in a 44-pin PCC and the SIO/3 is packaged in a 44-pin
QFP. Note that SIO/3 is only available in CMOS and in QFP package.}
\end{dscols}

\dsrule
\dshead{PIN DESCRIPTION}

\begin{dscols}
Figures 1 through 6 illustrate the three 40-pin configurations (bonding
options) available in the Z80 SIO (hereafter referred to as SIO or Z80 SIO).
The constraints of a 40-pin package make it impossible to bring out the
Receive Clock (\ol{RxC}), Transmit Clock (\ol{TxC}), Data Terminal Ready
(\ol{DTR}) and Sync (\ol{SYNC}) signals for both channels. Therefore, either
Channel B lacks a signal or two signals are bonded together.

\begin{featlist}
\item Z80 SIO/2 lacks \ol{SYNCB}
\item Z80 SIO/1 lacks \ol{DTRB}
\end{featlist}

\columnbreak

\begin{featlist}
\item Z80 SIO/0 has all four signals, but \ol{TxCB} and \ol{RxCB} are bonded
      together
\end{featlist}

\vspace{5pt}
\textbf{The 44-pin package, the Z80 SIO/4 for PLCC package, and Z80 SIO/3 for
QFP, has all options (Figure 7a and 7b).}

The first bonding option above (SIO/2) is the preferred version for most
applications. The pin descriptions are as follows:

\pin{B/\ol{A}.}{Channel A or B Select}{(input, High selects Channel B).
This input defines which channel is accessed during a data}
\end{dscols}

\pagefootrule
\clearpage

% ------------------------------- PAGE 100 ----------------------------------
% (style proof: Figures 1 and 2 only; Figures 3 and 4 follow in the full build)
\dsrule
\vspace{6pt}

\noindent
\begin{minipage}[t]{0.56\textwidth}
  \centering
  % Figure 1. Pin Functions -- Z80 SIO/2
% Geometry knobs: BW = body width, PY = row pitch.
\begin{tikzpicture}[x=1cm,y=1cm,baseline=(current bounding box.north)]
  \def\BW{2.30}          % chip body width
  \def\BH{7.60}          % chip body height
  \def\LX{-0.78}         % left  lead end
  \def\RX{3.08}          % right lead end
  \def\LB{-0.98}         % left  brace x
  \def\RBa{3.22}         % right inner brace x (modem control)
  \def\RBb{4.34}         % right outer brace x (channel)

  \draw[chipbox] (0,0) rectangle (\BW,-\BH);

  % ---- left side: name inside the body, lead + arrowhead outside ----------
  % #1 = y, #2 = label, #3 = direction (in|out|bi)
  \newcommand{\lsig}[3]{%
    \node[figfont, anchor=west] at (0.10,#1) {#2};
    \ifthenelse{\equal{#3}{in}}{\draw[sigarrow]  (\LX,#1) -- (0,#1);}{}%
    \ifthenelse{\equal{#3}{out}}{\draw[sigarrowr](\LX,#1) -- (0,#1);}{}%
    \ifthenelse{\equal{#3}{bi}}{\draw[sigboth]   (\LX,#1) -- (0,#1);}{}%
  }
  \newcommand{\rsig}[3]{%
    \node[figfont, anchor=east] at (\BW-0.10,#1) {#2};
    \ifthenelse{\equal{#3}{in}}{\draw[sigarrowr] (\BW,#1) -- (\RX,#1);}{}%
    \ifthenelse{\equal{#3}{out}}{\draw[sigarrow] (\BW,#1) -- (\RX,#1);}{}%
    \ifthenelse{\equal{#3}{bi}}{\draw[sigboth]   (\BW,#1) -- (\RX,#1);}{}%
  }

  % CPU data bus
  \lsig{-0.50}{\D{0}}{bi}   \lsig{-0.86}{\D{1}}{bi}
  \lsig{-1.22}{\D{2}}{bi}   \lsig{-1.58}{\D{3}}{bi}
  \lsig{-1.94}{\D{4}}{bi}   \lsig{-2.30}{\D{5}}{bi}
  \lsig{-2.66}{\D{6}}{bi}   \lsig{-3.02}{\D{7}}{bi}
  % Control from CPU
  \lsig{-3.62}{\ol{CE}}{in}    \lsig{-3.98}{\ol{RESET}}{in}
  \lsig{-4.34}{\ol{M1}}{in}    \lsig{-4.70}{\ol{IORQ}}{in}
  \lsig{-5.06}{\ol{RD}}{in}    \lsig{-5.42}{C/\ol{D}}{in}
  \lsig{-5.78}{B/\ol{A}}{in}
  % Daisy chain
  \lsig{-6.38}{\ol{INT}}{out}  \lsig{-6.74}{IEI}{in}
  \lsig{-7.10}{IEO}{out}

  % ---- right side ---------------------------------------------------------
  \rsig{-0.42}{RxDA}{in}        \rsig{-0.78}{\ol{RxCA}}{in}
  \rsig{-1.14}{TxDA}{out}       \rsig{-1.50}{\ol{TxCA}}{in}
  \rsig{-1.86}{\ol{SYNCA}}{bi}  \rsig{-2.22}{\ol{W/RDYA}}{out}
  \rsig{-2.66}{\ol{RTSA}}{out}  \rsig{-3.02}{\ol{CTSA}}{in}
  \rsig{-3.38}{\ol{DTRA}}{out}  \rsig{-3.74}{\ol{DCDA}}{in}

  \rsig{-4.30}{RxDB}{in}        \rsig{-4.66}{\ol{RxCB}}{in}
  \rsig{-5.02}{TxDB}{out}       \rsig{-5.38}{\ol{TxCB}}{in}
  \rsig{-5.74}{\ol{W/RDYB}}{out}
  \rsig{-6.18}{\ol{RTSB}}{out}  \rsig{-6.54}{\ol{CTSB}}{in}
  \rsig{-6.90}{\ol{DTRB}}{out}  \rsig{-7.26}{\ol{DCDB}}{in}

  % ---- device name --------------------------------------------------------
  \node[figfont, align=center, font=\sffamily\fontsize{5.4}{6}\selectfont]
    at (\BW/2,-3.32) {Z80\\SIO/2};

  % ---- power / clock ------------------------------------------------------
  \foreach \x/\lab in {0.55/+5\,V, 1.15/GND, 1.75/CLK} {
    \draw[sigarrow] (\x,-8.20) -- (\x,-\BH);
    \node[figfont, anchor=north] at (\x,-8.28) {\lab};
  }

  % ---- grouping braces ----------------------------------------------------
  % left braces bulge toward the labels (draw bottom-to-top)
  \draw[grpbrace] (\LB,-3.02) -- (\LB,-0.50);
  \node[figfont, anchor=east, align=center] at (\LB-0.14,-1.76) {CPU\\DATA BUS};

  \draw[grpbrace] (\LB,-5.78) -- (\LB,-3.62);
  \node[figfont, anchor=east, align=center] at (\LB-0.14,-4.70) {CONTROL\\FROM CPU};

  \draw[grpbrace] (\LB,-7.10) -- (\LB,-6.38);
  \node[figfont, anchor=east, align=center] at (\LB-0.14,-6.74)
    {DAISY\\CHAIN\\INTERRUPT\\CONTROL};

  % right braces bulge toward the labels (draw top-to-bottom)
  \draw[grpbrace] (\RBa,-2.66) -- (\RBa,-3.74);
  \node[figfont, anchor=west, align=center] at (\RBa+0.14,-3.20) {MODEM\\CONTROL};
  \draw[grpbrace] (\RBa,-6.18) -- (\RBa,-7.26);
  \node[figfont, anchor=west, align=center] at (\RBa+0.14,-6.72) {MODEM\\CONTROL};

  \draw[grpbrace] (\RBb,-0.42) -- (\RBb,-3.74);
  \node[figfont, anchor=west] at (\RBb+0.14,-2.08) {CHANNEL A};
  \draw[grpbrace] (\RBb,-4.30) -- (\RBb,-7.26);
  \node[figfont, anchor=west] at (\RBb+0.14,-5.78) {CHANNEL B};
\end{tikzpicture}

  \figcap{1}{Pin Functions}
\end{minipage}%
\hfill
\begin{minipage}[t]{0.42\textwidth}
  \centering
  % Figure 2. 40-pin Dual-In-Line Package (DIP), Pin Assignments -- Z80 SIO/2
% Pin lists read top-to-bottom: \leftpins is pins 1..20, \rightpins is 40..21.
\begin{tikzpicture}[x=1cm,y=1cm,baseline=(current bounding box.north)]
  \def\BW{1.52}      % package body width
  \def\PY{0.335}     % pin pitch
  \def\PT{0.45}      % y of pin 1
  \def\SW{0.155}     % pin stub width
  \def\SH{0.085}     % pin stub half-height

  \pgfmathsetmacro{\BH}{\PT+19*\PY+0.45}
  \draw[chipbox] (0,0) rectangle (\BW,-\BH);
  % index notch
  \draw[line width=0.4pt] (\BW/2-0.16,0) arc (180:360:0.16);

  % left column: pins 1..20
  \foreach \n/\lab [count=\i from 0] in {%
    1/\D{1}, 2/\D{3}, 3/\D{5}, 4/\D{7}, 5/\ol{INT}, 6/IEI, 7/IEO, 8/\ol{M1},
    9/{+5\,V}, 10/\ol{W/RDYA}, 11/\ol{SYNCA}, 12/RxDA, 13/\ol{RxCA}, 14/\ol{TxCA},
    15/TxDA, 16/\ol{DTRA}, 17/\ol{RTSA}, 18/\ol{CTSA}, 19/\ol{DCDA}, 20/CLK} {
    \pgfmathsetmacro{\y}{-\PT-\i*\PY}
    \draw[line width=0.35pt] (-\SW,\y-\SH) rectangle (0,\y+\SH);
    \node[figfont, anchor=west] at (0.09,\y) {\n};
    \node[figfont, anchor=east] at (-\SW-0.08,\y) {\lab};
  }

  % right column: pins 40..21
  \foreach \n/\lab [count=\i from 0] in {%
    40/\D{0}, 39/\D{2}, 38/\D{4}, 37/\D{6}, 36/\ol{IORQ}, 35/\ol{CE}, 34/{B/\ol{A}},
    33/{C/\ol{D}}, 32/\ol{RD}, 31/GND, 30/\ol{W/RDYB}, 29/RxDB, 28/\ol{RxCB},
    27/\ol{TxCB}, 26/TxDB, 25/\ol{DTRB}, 24/\ol{RTSB}, 23/\ol{CTSB},
    22/\ol{DCDB}, 21/\ol{RESET}} {
    \pgfmathsetmacro{\y}{-\PT-\i*\PY}
    \draw[line width=0.35pt] (\BW,\y-\SH) rectangle (\BW+\SW,\y+\SH);
    \node[figfont, anchor=east] at (\BW-0.09,\y) {\n};
    \node[figfont, anchor=west] at (\BW+\SW+0.08,\y) {\lab};
  }

  \node[figfont, align=center, font=\sffamily\fontsize{5.4}{6}\selectfont]
    at (\BW/2,-3.55) {Z80\\SIO/2};
\end{tikzpicture}

  \figcap{2}{40-pin Dual-In-Line Package (DIP),\\Pin Assignments}
\end{minipage}

\vspace{8pt}
\dsrule

\pagefootrule
\end{document}
